\chapter{Аудит геометрических кандидатов на абсолютную энергию часов}
\label{ch:v10-cell-birth-clock-energy-geometric-anchor-audit}

\section{Контракт допустимого якоря}

Предыдущая глава вывела относительное тождество
\begin{equation}
 \frac{H_B}{\Omega}=\Delta\zeta,
 \qquad \Omega=\frac{E_C}{\hbar},
 \label{eq:v10-anchor-relative-growth}
\end{equation}
но оставила одномерную орбиту общего масштаба. Геометрический кандидат на
$E_C$ допустим только при одновременном выполнении пяти условий:
\begin{enumerate}
\item он имеет размерность энергии;
\item существует внутри уже построенного носителя;
\item не определяется через искомую частоту часов;
\item имеет явное отображение в коэффициент $E_C H_0$;
\item разрывает масштабную орбиту без наблюдаемой цели или внешней константы.
\end{enumerate}

\section{Реестр кандидатов}

Проверены девять классов:
\begin{equation}
 \left\{
 \frac{\hbar}{\tau_{\rm cell}},
 \frac{\hbar c}{\ell_{\rm cell}},
 \hbar c\sqrt{\frac{\Lambda_{\rm growth}}3},
 \hbar c\Lambda_{\rm sp},
 \frac{\hbar c}{24R},
 k_BT,
 \sqrt{\frac{\hbar c^5}{G}},
 \hbar H_{\rm obs},
 S_{\rm vac}
 \right\}.
 \label{eq:v10-anchor-candidate-menu}
\end{equation}
Индикаторная матрица девяти строк и пяти условий имеет ранг $5$, то есть
реестр действительно различает все требования. Однако произведение каждой
строки равно нулю:
\begin{equation}
 \mathbf p=(0,0,0,0,0,0,0,0,0)^T.
 \label{eq:v10-anchor-pass-vector}
\end{equation}
Наибольший счёт равен лишь $3/5$. Ни один внутренне доступный кандидат не
разрывает общую масштабную орбиту.

Собственные время и длина ячейки имели бы правильный тип, но их величины
ещё не выбраны:
\begin{equation}
 E_\tau\tau_{\rm cell}=\hbar,
 \qquad
 E_\ell\ell_{\rm cell}=\hbar c.
 \label{eq:v10-anchor-cell-products}
\end{equation}
Эти равенства фиксируют произведения, а не отдельные множители. Аналогично
казимировский кандидат сохраняет свободный радиус:
\begin{equation}
 24E_{\rm Cas}R=\hbar c.
 \label{eq:v10-anchor-casimir-product}
\end{equation}
Спектральное обрезание и температура КМС также не имеют внутреннего
селектора абсолютной величины. Планковская энергия и наблюдаемый темп
Хаббла разорвали бы орбиту, но лишь после внесения $G$ или наблюдаемой цели.

\section{Круговость космологического кандидата}

Особого внимания требует гипотеза, что энергия часов задаётся
космологической постоянной. В текущей ветви сама величина роста определена
через те же часы:
\begin{equation}
 H_B=\Delta\zeta\frac{E_C}{\hbar},
 \qquad
 \Lambda_{\rm growth}=3\left(\frac{H_B}{c}\right)^2.
 \label{eq:v10-anchor-growth-curvature}
\end{equation}
Поэтому естественная энергия кривизны равна
\begin{equation}
 E_\Lambda
 :=\hbar c\sqrt{\frac{\Lambda_{\rm growth}}3}
 =\Delta\zeta E_C.
 \label{eq:v10-anchor-curvature-energy}
\end{equation}
Попытка восстановить энергию часов даёт
\begin{equation}
 \frac{E_\Lambda}{\Delta\zeta}=E_C,
 \label{eq:v10-anchor-curvature-tautology}
\end{equation}
то есть точное тождество, а не новое уравнение отбора. Космологическая
постоянная может стать независимым якорем только если она выведена из
другого геометрического сектора, не использующего $H_B$ или $E_C$.

\section{Линейный масштабный запрет}

В логарифмических координатах $(E_C,\Omega,\Gamma_B,H_B)$ три уже
установленные относительные связи представлены матрицей
\begin{equation}
 A_{\rm rel}=
 \begin{pmatrix}
 -1&1&0&0\\
 0&-1&1&0\\
 0&0&-1&1
 \end{pmatrix},
 \qquad \rank A_{\rm rel}=3,
 \qquad \dim\ker A_{\rm rel}=1.
 \label{eq:v10-anchor-relative-map}
\end{equation}
Вектор $(1,1,1,1)^T$ лежит в ядре. Дополнительная строка, фиксирующая
любую одну абсолютную энергию, повысила бы ранг до $4$, но среди внутренних
кандидатов такой строки нет.

\begin{theorem}[Исчерпание текущих геометрических кандидатов]
Среди собственного времени и длины ячейки, кривизны роста, спектрального
обрезания, казимировской энергии, температуры КМС, планковской энергии,
наблюдаемого темпа Хаббла и $S_{\rm vac}$ нет внутреннего независимого
якоря $E_C$. Реестр покрывает все пять условий, но даёт $0/9$ проходов.
\end{theorem}

\begin{nogocriterion}[Запрет кругового масштаба через рост]
Нельзя использовать $\Lambda_{\rm growth}=3(H_B/c)^2$ для вывода той же
энергии часов, через которую определён $H_B$. Получаемое равенство
$E_\Lambda/\Delta\zeta=E_C$ тождественно и не снимает масштабную орбиту.
\end{nogocriterion}

\begin{protocol}[Статус аудита геометрического якоря]\begin{enumerate}
\item проверенных классов кандидатов: \textbf{$9$};
\item условий физического происхождения: \textbf{$5$};
\item ранг матрицы условий: \textbf{$5$};
\item прошедших кандидатов: \textbf{$0/9$};
\item внутренних разрывателей орбиты: \textbf{$0$};
\item энергия часов: \textbf{$0/1$};
\item физический космологический масштаб: \textbf{$0/1$};
\item следующий гейт: \textbf{родитель собственного четырёхмерного объёма
ячейки}.
\end{enumerate}\end{protocol}

Машинный сертификат сохранён в
\path{s2t_v10_cell_birth_clock_energy_geometric_anchor_candidate_audit_gate_results.json}.